本报告完成了《面向对象程序设计实践》选题 56 ——“3D 点云目标分割方法设计与实现”。
项目用 C++17 从零实现了一套点云目标分割系统：先用 KD 树组织点云并做 $k$ 近邻查询，
再用主成分分析（PCA）估计每个点的法向量与曲率，然后以\textbf{区域生长}算法按表面朝向
的一致性把点云划分成若干光滑片段，并通过逐层收紧法向夹角阈值实现\textbf{由粗到细}的
迭代细分。系统在 GLFW + OpenGL + Dear ImGui 上构建了三维交互界面，支持点云的旋转、
缩放、平移，分割结果按片段分色显示，并能用滑块在各细分层之间回放由粗到细的分割过程。

工程采用面向对象思想组织：把向量、点云、KD 树、法向量估计、分割算法、文件读写等分别
封装为独立的类与模块，核心算法库与图形界面完全解耦，从而既能在命令行下批处理、也能在
界面中交互。项目自带合成点云生成器并配套 69 项单元测试，覆盖 KD 树近邻、平面法向量、
PLY 读写以及分割正确性，全部通过；在 9000 点的演示场景上，法向量估计与三层分割合计
约 80\,毫秒完成，最细层对“地面 + 球 + 圆柱 + 斜面”四类目标的分割纯度达到 95\% 以上。
报告还讨论了纯区域生长在锐利棱边处的“泄漏”局限及其成因。
